\documentclass[journal]{IEEEtran}
\usepackage{amsmath, amssymb}
\usepackage{graphicx}
\usepackage{booktabs}
\usepackage{hyperref}

\title{An IEEE Journal Paper Template}
\author{First~Author,~Second~Author,~and~Third~Author%
\thanks{Manuscript received ~; This work was supported by ~.%
\newline \textit{First Author} is with the Department of Computer Science,
Example University (e-mail: first@example.edu).}}

\begin{document}
\maketitle

\begin{abstract}
An abstract for the IEEE Transactions style. Replace with a concise summary
of your contribution, methods, and results in roughly 150--250 words.
\end{abstract}

\begin{IEEEkeywords}
LaTeX, IEEEtran, journal, template
\end{IEEEkeywords}

\section{Introduction}
Journal papers demand a thorough literature review and an in-depth
evaluation. This template provides the IEEE Transactions two-column
layout~\cite{ref1}.

\section{Related Work}
Discuss prior art and position your contribution~\cite{ref2}.

\section{Proposed Method}
Define the optimization problem:
\begin{equation}
  \min_{\theta} \; \mathbb{E}_{(x,y)\sim\mathcal{D}}
  \big[\ell(f_\theta(x), y)\big] + \lambda\Omega(\theta).
\end{equation}

\section{Experiments}
\begin{table}[!t]
  \centering
  \caption{Example results.}
  \label{tab:r}
  \begin{tabular}{lcc}
    \toprule
    Method & Metric-1 & Metric-2 \\
    \midrule
    Baseline & 81.2 & 0.79 \\
    Ours     & \textbf{88.1} & \textbf{0.87} \\
    \bottomrule
  \end{tabular}
\end{table}
Table~\ref{tab:r} summarizes the comparison.

\section{Conclusion}
Summarize findings and future work.

\begin{thebibliography}{9}
\bibitem{ref1} IEEE Editorial Staff, ``Author guidelines,'' \emph{IEEE Trans.}, 2024.
\bibitem{ref2} A. Author, B. Author, ``A related article,'' in \emph{Proc. Conf.}, 2023.
\end{thebibliography}

\end{document}
